\documentclass[../main-en.tex]{subfiles}
\begin{document}
\section{Introduction}

The Sensor Data Processing Hub collects readings from environmental sensors (temperature, CO$_2$, humidity), applies a chain of transformations to them (unit conversion, smoothing, threshold alerts), and sends the results to a console dashboard and a log file. The sensors are simulated; we implement two sources, a CSV file reader and a random generator, and treat a socket feed as a future change (\cref{sec:scenario-socket}). The unit conversion and the transformation chain are defined in a configuration file.

This document describes the software architecture of the hub. The approach follows the IKT300 Software Architecture and Design course~\cite{wiklund_ikt300_2026}, based on Software Architecture in Practice~\cite{bass_software_2021}: architecture is important because it enables communication among stakeholders and its early decisions constrain the quality attributes of the system~\cite{wiklund_ch2_nodate}. Our main quality attribute is modifiability, since new sources, transformations and sinks are expected. The architecture is a central mediator, \texttt{Hub}, that pulls readings from the sources, passes them through a pipe-and-filter chain of transformations and pushes the results to the sinks.

The report presents the module view (\cref{sec:module-view}), the component-and-connector view (\cref{sec:cc-view}) and the allocation view (\cref{sec:allocation-view}). \Cref{sec:mapping} maps elements between the views, \cref{sec:rationale} justifies the key architectural decisions, and \cref{sec:modifiability} states modifiability scenarios with concrete response measures. The C\# code and a short README that explains how to build and run it are delivered with this report.

\end{document}
